%%% PART
% Macro pour souligner le texte (en gras et rouge)
\newcommand{\underlinetitlePart}[1]{%
  \tikz[baseline=(text.base)]{
    \node[inner sep=0pt, text=black] (text) {#1};
    \draw[line width=2pt,background2] 
      ($(text.south west)-(0,4pt)$) -- ($(text.south east)-(0,4pt)$);
  }%
}

% Macro pour \part avec rectangle noir collé à droite
\newcommand{\parttitlewithbox}[1]{%
 https://github.com/jeanollion/bacmman/wiki \behttps://github.com/jeanollion/bacmman/wikigin{tikzpicture}
    % boîte de largeur \textwidth
    \node[inner sep=0pt, anchor=west] (text) at (0,0) {%
      \hspace{0cm}\underlinetitlePart{{\color{black}\MakeUppercase{#1}}}%
    };
    % rectangle collé à droite (bord de la page)
    \filldraw[primarycolor] 
      ($(0.9\textwidth,0) + (0,-0.5em)$) rectangle ++(0.5em,1em);
  \end{tikzpicture}%
}

% Formatage de \part
\titleformat{\part}[block]
  {\robotoCondensed\bfseriesfontsize{26}{30}\selectfont}
  {\newpage}
  {0pt}
  {\parttitlewithbox}
\titlespacing*{\part}{1cm}{0cm}{0cm}